\subsection{Problem statement and control objectives}
\textcolor{red}{SHOULD WE CANCEL THIS PART?}
This report details the design, simulation, and experimental validation of control strategies for a magnetic levitation system.
Controlling this system presents two fundamental physical challenges: it is inherently open-loop unstable (the ball cannot naturally hover),
and the electromagnetic attractive force is highly non-linear.
Consequently, standard linear control formulations are only strictly valid within a bounded region around a defined equilibrium point. 

Control Objectives To safely stabilize the plant and overcome these non-linearities, the control architecture is decoupled into this progressive objectives: 

\begin{enumerate}[label=\textbf{O\arabic*.},leftmargin=1.4cm]
    \item \textbf{Coil current control.} The electrical dynamics of the coil are regulated by an inner current loop.
    This loop is designed to be faster than the mechanical dynamics, so that the outer position controller can command a current reference instead of directly commanding the coil voltage.
    \item \textbf{Ball position stabilization.} The unstable mechanical dynamics are stabilized around a nominal levitation height. Frequency-based and state-space techniques are designed, simulated and compared.
\end{enumerate}
